\section{Discussion: challenges, validation and outlook}\label{sec:discussion}

The Self-Aware Steel Structure framework proposed in this work is conceptually coherent and grounded in mature theory, yet its transition from paradigm to deployable practice raises a set of interdependent challenges that must be addressed candidly. The most fundamental caveat is that the framework is, at present, a proposal: the testbed, sensing network, digital twin and AI layer described above have not yet been built or experimentally validated. The arguments below should therefore be read as a critical roadmap rather than as evidence of demonstrated capability.

\textbf{Validation strategy.} A credible validation must proceed in stages and must distinguish between simulated and physically induced damage. The FEM provides a controllable environment in which stiffness reductions, connection releases and mass changes can be parametrically introduced to generate a hybrid dataset of healthy and damaged responses for initial AI training, a practice with clear precedent in using a validated finite element model as a data source \citep{naser2023ml}. FEM-simulated damage inherits the limitation that ``FEA is simulation, not reality'' \citep{cook2002fea}: even an accurate analysis of an inadequate mathematical model can disagree with the physical structure. Models trained predominantly on synthetic cases may therefore learn artefacts of the idealization, such as the particular connection-stiffness model adopted, rather than true structural physics, and real or experimental data are generally to be favored over purely synthetic data \citep{naser2023ml}. A classifier assessed only on cases that the finite-element model itself generated and labelled demonstrates class separability under the model's own assumptions, not performance on the physical structure; any preliminary numerical study is therefore reported as a specification of the validation protocol and the target metrics it will use, rather than as evidence of field capability. Physically induced damage, through controlled bolt loosening, added masses, or member removal in the bracing, is indispensable for closing this reality gap. Bolted-connection loosening is a particularly appropriate target because it alters system connectivity and adds energy dissipation, which typically manifests as a measurable increase in damping \citep{farrar2013shm,razi2013bolted,chelimilla2023looseness}, consistent with the documented precedent in which period lengthening accompanied stiffness reduction and increased damping at damaging amplitudes \citep{chopra2017dynamics}. The contrast between this self-interpreting workflow and conventional data-collecting monitoring is summarized in Table~\ref{tab:conv-vs-aware}. Validation must be grouped rather than randomly shuffled: folds must hold out whole \textit{tests}, damage \textit{scenarios} and \textit{operating conditions} together, rather than individual and possibly adjacent windows, while the class imbalance between abundant healthy and scarce damaged data is handled explicitly, to guard against the optimistic bias and overfitting that arise when training performance far exceeds testing performance \citep{naser2023ml,figueiredo2010machine}.

\textbf{Generalization across structures.} A model trained on one SHS braced frame encodes that structure's specific mass, stiffness and modal signature. Because natural frequencies and mode shapes depend only on mass and stiffness \citep{chopra2017dynamics}, they are intrinsically structure-specific, and a damage classifier learned on the testbed will not transfer directly to a differently configured frame. Generalization will likely require physics-informed features, such as modal parameters and prediction residuals against the structure's own digital twin, rather than raw signals, together with transfer-learning strategies \citep{naser2023ml,karaaslan2021attention} that adapt a pretrained model to a new structure using limited local data.

\textbf{Sensor reliability, cost and computation.} The distributed network of strain gauges, accelerometers, temperature, displacement and inclination sensors introduces an inherent cost-versus-fidelity trade-off in deployment \citep{farrar2013shm}. Sensors degrade, drift and fail, and ``sensors cannot measure damage'' directly: feature extraction and statistical classification are always required \citep{farrar2013shm}. Numerically driven sensor placement, supported by reduced-order modeling and sparse or optimal placement techniques \citep{brunton2019data,tan2020sensorplacement,hassani2023sensorplacement}, can reduce channel count while preserving observability, but introduces its own approximation. Real-time operation is constrained by the cost of repeatedly solving the eigenproblem and the harmonic and transient analyses, which is far greater than the static solution \citep{cook2002fea}; reduced-order models and data-driven modal decomposition (DMD) \citep{brunton2019data} are the natural mitigations, at the price of fidelity in higher modes.

\textbf{Uncertainty, interpretability and trust.} A self-aware structure that \textit{reports} its state must also report its confidence. Without data normalization separating benign operational and environmental variability from genuine damage, the system will produce false-positive indications \citep{farrar2013shm,figueiredo2010machine}. Temperature-induced stiffness change is the canonical confounder, and the thermal-compensation channel addresses it directly. It is not the only one, however: varying operational loads, wind and humidity also perturb the modal parameters, and since the testbed does not instrument these quantities, they have to be separated statistically rather than measured away. Conditioning the healthy model on the inferred operational state, instead of assuming that temperature exhausts the benign variability, is what keeps such effects from being read as damage. The axiomatic trade-off between damage sensitivity and noise rejection \citep{farrar2013shm} bounds what any algorithm can achieve. Engineer trust further demands interpretability: black-box classifiers should be accompanied by explainability tools such as permutation feature importance, SHAP and partial dependence plots, so that flagged anomalies map to physically checkable indicators \citep{naser2023ml}. Deep models also face generalizability and interpretability limits together with large-data requirements \citep{brunton2019data}.

\textbf{Standardization, integration and scalability.} For the paradigm to reach asset-management practice it must integrate with BIM and digital-asset workflows and adopt standardized data and damage-state semantics, for example the Rytter hierarchy of existence, location, type, severity and prognosis \citep{farrar2013shm}, so that autonomous outputs are interpretable across stakeholders \citep{chiachio2022structural,mariniello2025enhancing}. Scaling from a laboratory frame to in-service warehouses, platforms and modular structures multiplies environmental variability, connection count and uncontrolled excitation, where the rotating-unbalance source is replaced by ambient and operational loads \citep{he2022integrated}.

\textbf{Multimodal perception.} The framework developed here is deliberately vibration-based: perception, the digital twin and the cognition layer all operate on the dynamic response of the SHS frame, and the testbed carries no optical instrumentation. The five-layer architecture is nonetheless open at its sensing and cognition layers to complementary perception modalities. Vision-based deep learning is the most immediate of these: on infrastructure imagery, convolutional networks already detect and classify surface damage such as corrosion and cracking \citep{karaaslan2021attention,arya2021roaddamage}, and segmentation variants can localize such damage at the pixel level. A visual channel of this kind would observe surface and geometric damage that modal features reach only indirectly, and fusing a learned visual damage index with the modal and twin-residual features could corroborate an event across independent modalities. Bolted-connection loosening, for instance, would register as both an increase in damping and a localized change in the image of the joint. This multimodal route is left to future work: it requires its own camera network, annotated image datasets and trained models, none of which belongs to the present vibration-based specification.

\textbf{Limitations.} Beyond the absence of experimental validation, the framework relies on a single excitation mechanism and on modally observable damage types; and although the bolted connections are now represented as semi-rigid rotational springs rather than ideal pins, their stiffness still has to be identified experimentally before the numerical model can serve as a trusted baseline. The detectable damage size is inversely proportional to the excitation frequency range \citep{farrar2013shm}, which limits sensitivity to small or local flaws. These constraints define the agenda for the experimental program rather than negating the concept.
